% ============================================================================
%  MULTICLASS (K>=3) knowability--capacity.  (theory_v2, June 2026)
%
%  Addresses the OPEN part (b) of Conjecture conj:gen-capacity in
%  paper/sections/knowability_capacity_general.tex: "the MULTICLASS case (K>=3),
%  where the 1-D location-model intuition breaks."
%
%  Status labels (as in knowability_capacity_general.tex):
%     (PROVEN)            -- complete elementary proof, no extra hypothesis.
%     (PROVEN, COND.)     -- proven UNDER an explicitly stated structural condition.
%     (PROVEN, WITNESS)   -- proven by an exhibited explicit witness/counterexample.
%     (OPEN)              -- precise conjecture with the exact missing piece.
%
%  HONEST HEADLINE (see Remark rmk:mc-scope): this is a STRONG PARTIAL on the
%  multiclass frontier, NOT a closure. It (i) gives the exact vector benefit identity
%  for K>=3; (ii) proves a single-conflict-pair multiclass capacity at tau=1 under an
%  explicit pair of hypotheses (the multiclass R1/R2 analogue); and (iii) proves TWO
%  necessity results that delimit exactly where the binary capacity stops: the
%  conflict-pair structure (M1) is necessary, and the standard binary observable margin
%  is a strictly-too-coarse reduct (third-class confounding). The genuinely
%  multi-conflict-pair / vector-concept regime remains OPEN.
%
%  Macros assumed (as in knowability_capacity_general.tex): \Delta \R \TV \sign \Phi
%  \mathbb E, environments theorem/lemma/remark, \boxed.
%  All numbers validated by theory_v2/val_multiclass_capacity.py (seed 20260629)
%  -> theory_v2/results_multiclass_capacity.json.
% ============================================================================
\section{The multiclass knowability--capacity: how far it lifts, and the exact wall}
\label{sec:mc-capacity}

% local macros so this section compiles standalone (mirror the companion sections)
\providecommand{\TV}{\mathrm{TV}}
\providecommand{\sgn}{\operatorname{sign}}
\providecommand{\R}{\mathbb R}

The general-capacity section (\S\ref{sec:knowcap-gen}) and the dichotomy
(\S\ref{sec:knowdich}) settle the \emph{binary} sign-of-benefit problem, including the
removal of the regularity conditions R1/R2 at the level of \emph{existence}
(Theorem~\ref{thm:dich-main}: a capacity exists iff the benefit sign factors through the
observable reduct). They leave open the second half of
Conjecture~\ref{conj:gen-capacity}: the \emph{multiclass} case $K\ge3$, ``where the
$1$-D location-model intuition breaks.'' This section pins down exactly how it breaks. The
answer is a sharp map: the binary capacity lifts \emph{verbatim} to a single
\emph{conflict pair} of classes, and stops there for two independent and explicit
reasons.

% ----------------------------------------------------------------------------
\subsection{The multiclass benefit and the new degree of freedom}\label{sec:mc-setup}
% ----------------------------------------------------------------------------
\paragraph{Model.} Covariate shift $Q_X=q_0(\cdot-\mu)$ on $\R$ (one dimension already
exposes the obstruction), with $q_0$ a fixed base density and $\mu$ the unknown shift,
estimable label-free from an unlabeled target sample. Labels $Y\in\{1,\dots,K\}$, $0/1$
loss. The deployed classifier $g$ and the candidate $f$ are measurable maps
$\mathcal X\to\{1,\dots,K\}$ with disagreement region $D=\{x:g(x)\neq f(x)\}$. The
\emph{concept} (target label channel) is $\eta(x)=(\eta_1(x),\dots,\eta_K(x))$ on the
probability simplex $\Delta^{K-1}$; under covariate shift $\eta$ does \emph{not} affect
$Q_X$. The deployment benefit is $\Delta=R_Q(g)-R_Q(f)$, and $\Delta>0\iff$ switching
$g\to f$ helps.

\begin{lemma}[Exact multiclass benefit identity; PROVEN]\label{lem:mc-identity}
On $D$ write $j_g(x)$ for the class predicted by $g$ and $j_f(x)$ for the class predicted
by $f$ (so $j_g(x)\neq j_f(x)$). Then the $0/1$ excess loss of $g$ over $f$ vanishes off
$D$ and equals $\mathbf 1[Y{=}j_f]-\mathbf 1[Y{=}j_g]$ on $D$, so
\begin{equation}\label{eq:mc-delta}
  \boxed{\;\Delta \;=\; \int_D \big(\eta_{j_f(x)}(x)-\eta_{j_g(x)}(x)\big)\,\mathrm dQ(x)
  \;=:\;\int_D \Lambda(x)\,\mathrm dQ(x),\;}
  \qquad \Lambda(x)\in[-1,1].
\end{equation}
The \emph{benefit field} $\Lambda(x)=\eta_{j_f}(x)-\eta_{j_g}(x)$ involves only the two
disagreeing coordinates. The residual mass
$m_{\mathrm{rest}}(x)=1-\eta_{j_g}(x)-\eta_{j_f}(x)\ge0$ on the other $K-2$ classes is a
\emph{free function}, invisible to $Q_X$.
\end{lemma}
\begin{proof}
Off $D$, $g=f$ and the losses coincide. On $D$, $g$ and $f$ are the deterministic labels
$j_g,j_f$, so $\mathbf 1[g\neq Y]-\mathbf 1[f\neq Y]=(1-\mathbf 1[Y{=}j_g])-(1-\mathbf
1[Y{=}j_f])=\mathbf 1[Y{=}j_f]-\mathbf 1[Y{=}j_g]$, with conditional mean
$\eta_{j_f}(x)-\eta_{j_g}(x)$. Integrating against $Q$ over $D$ gives \eqref{eq:mc-delta};
$|\Lambda|\le1$ since $\eta_{j_f},\eta_{j_g}\ge0$ and $\eta_{j_f}+\eta_{j_g}\le1$.
\emph{(Validated against a $4\times10^6$-sample Monte-Carlo of $R_Q(g)-R_Q(f)$, $K=3$:
$|\Delta_{\mathrm{MC}}-\Delta_{\mathrm{c.f.}}|\le3.6\times10^{-4}$;
\texttt{val\_multiclass\_capacity.py} Block~0.)}
\end{proof}

\noindent\emph{Why this is genuinely new.} In binary the disagreeing pair is forced to be
$\{0,1\}$ and $\Lambda=\eta_1-\eta_0=2\eta_1-1$ is a single-coordinate scalar pinned by
the (scalar) concept, so $\Lambda$ inherits the single-crossing/monotone structure of a
one-dimensional threshold. For $K\ge3$, $\Lambda$ is a \emph{difference of two simplex
coordinates}, and the residual $m_{\mathrm{rest}}$ is an extra, unobservable degree of
freedom. Both facts have consequences, isolated below.

% ----------------------------------------------------------------------------
\subsection{The single-conflict-pair extension: the capacity lifts}\label{sec:mc-lift}
% ----------------------------------------------------------------------------
We make explicit the multiclass analogue of R1 (unique flip locus) and R2
(monotone-in-nuisance).
\begin{itemize}\itemsep2pt
\item[\textnormal{(M1)}] \textbf{Single conflict pair.} There are \emph{fixed} labels
$j_g\neq j_f$ such that $g\equiv j_g$ and $f\equiv j_f$ on \emph{all} of $D$. (Automatic in
binary; holds whenever $g,f$ differ along a single decision boundary.)
\item[\textnormal{(M2)}] \textbf{MLR location nuisance.} $q_0$ is log-concave (MLR in the
location), and the source contrast $\Lambda_S$ has a monotone, single-crossing profile in
the sufficient direction, so the flip locus $\theta_0(\mu)$ is a monotone, observable
function of the unknown shift (exactly the hypothesis of Theorem~\ref{thm:gen-mlr}).
\item[\textnormal{(M3)}] \textbf{Contrast-mass drift budget $\varepsilon$.} The admissible
target concept's single crossing of $\Lambda$ moves by at most $\varepsilon$ in the
observable mass coordinate $u=\Phi_{q_0}(\theta-\mu)$ of that crossing (the multiclass
analogue of the boundary-mass budget \eqref{eq:knowcap-class}, now placed on the
\emph{contrast} field rather than on a top score).
\end{itemize}

\begin{theorem}[Single-conflict-pair multiclass capacity; PROVEN, COND.]
\label{thm:mc-singlepair}
Assume \textnormal{(M1)--(M3)} and $Q_X$ known. Then the multiclass benefit
\eqref{eq:mc-delta} reduces \emph{exactly} to the binary single-crossing benefit in the
scalar field $\Lambda$, and the population knowability--capacity
(Theorem~\ref{thm:knowcap}) holds verbatim with $\Phi\to\Phi_{q_0}$ and the boundary mass
read off $\Lambda$'s crossing: with $\tau=1$,
\[
  K \;=\; \frac{\big|\,u_S-u_0\,\big|}{\varepsilon},\qquad
  u_S=\Phi_{q_0}(\theta_S-\mu),\quad
  u_0=\text{band-median mass of }D\text{ along }\Lambda,
\]
is a \emph{capacity}: $K>1\iff\sgn\Delta$ is label-free identifiable over the admissible
class, with the plug-in certificate $\sgn(u_S-u_0)$ correct with error $0$; and $K<1$
gives two admissible concepts with the same $Q_X$ and opposite $\sgn\Delta$, hence
minimax label-free error $\tfrac12$. The finite-$n$ Le~Cam boundary layer
(Prop.~\ref{prop:knowcap-finiten}) carries over with $\mu$ estimated label-free at rate
$O(1/\sqrt n)$.
\end{theorem}
\begin{proof}
Under (M1) the labels $j_g,j_f$ are constant on $D$, so $\Lambda(x)=\eta_{j_f}(x)-\eta_{j_g}(x)$
is a single scalar field and the residual $m_{\mathrm{rest}}(x)$ enters \eqref{eq:mc-delta}
\emph{only} through $\eta_{j_g}+\eta_{j_f}=1-m_{\mathrm{rest}}$, which is orthogonal to the
contrast: writing $\eta_{j_f}-\eta_{j_g}=(1-m_{\mathrm{rest}})\,c(x)$ with
$c(x)\in[-1,1]$ the within-pair contrast, $\Delta=\int_D(1-m_{\mathrm{rest}})\,c\,\mathrm dQ$.
With $m_{\mathrm{rest}}$ a constant reservoir $c_3$ on $D$ (the canonical (M1) case), this
is $(1-c_3)$ times the binary benefit of Lemma~\ref{lem:knowcap-identity} for the threshold
concept $c$; the factor $(1-c_3)>0$ does not affect $\sgn\Delta$. By (M2) the flip locus is
the band median along $\Lambda$, monotone and observable in $\mu$; by (M3) the admissible
set is the mass interval $[u_S-\varepsilon,u_S+\varepsilon]$. Theorem~\ref{thm:knowcap}
applies verbatim: $K>1$ iff $u_0$ is outside the admissible interval (achievability),
$K<1$ iff inside (converse, two concepts sharing $Q_X$ with opposite sign, $\TV=0$). The
finite-$n$ layer is Proposition~\ref{prop:knowcap-finiten} with $\Phi\to\Phi_{q_0}$.
\emph{(Validated: $50{,}000$ random draws with the residual reservoir
$c_3\in\{0.0,\dots,0.6\}$, equivalence $K>1\iff$ identifiable with $0$ mismatches off a
$10^{-6}$ shell, seed-robust; \texttt{val\_multiclass\_capacity.py} Block~A.)}
\end{proof}

\noindent This is a genuine multiclass extension: it is \emph{not} reducible to ``ignore
the extra classes,'' because the residual mass is present and unobservable; what (M1)
buys is that the residual cancels in the contrast, restoring a scalar single-crossing
problem.

% ----------------------------------------------------------------------------
\subsection{Necessity I: the conflict-pair structure (M1) cannot be dropped}
\label{sec:mc-breakerM1}
% ----------------------------------------------------------------------------
The binary case lifts to MLR/log-concave \emph{per-class} score families
(Theorem~\ref{thm:gen-mlr}); one might hope the same regularity suffices in multiclass.
It does not, and the obstruction is structural, not a pathology of the base density.

\begin{theorem}[Multi-conflict-pair breaker; PROVEN, WITNESS]\label{thm:mc-breakerM1}
There are canonical multiclass classifiers $g,f$ (argmax of affine per-class scores --- the
exact per-coordinate MLR/log-concave regularity that suffices in binary) for which (M1)
fails: $D$ decomposes into sub-bands carrying \emph{different} conflict pairs
$(j_g,j_f)$. Then the benefit field $\Lambda$ of \eqref{eq:mc-delta} is piecewise (a
different coordinate pair on each sub-band), its sign crosses several times in $x$, and the
benefit $\Delta(\mu)=\int_D\Lambda\,\mathrm dQ_\mu$ \textbf{oscillates} in the unknown shift
$\mu$. The label-free identifiable set $\{\mu:\sgn\Delta\text{ identifiable}\}$ therefore
\textbf{fragments into $\ge3$ components}, while the binary single-threshold baseline stays
monotone (one flip). Consequently no monotone (binary-R2) flip locus and no single scalar
threshold $\tau$ in any margin built along one direction governs identifiability: (M1) is
\emph{necessary} for the scalar capacity, and the binary regularity does not lift.
\end{theorem}
\begin{proof}[Proof (by exhibited witness)]
Take $D$ split into three sub-bands on which the argmax disagreement pairs give within-pair
contrasts of signs $+,-,+$ (a ``sandwich''; realizable by two affine-score argmax
classifiers whose decision boundaries interleave). The resulting field $\Lambda$ has
$|\Lambda|\le1$ (Lemma~\ref{lem:mc-identity}, so it is a bona fide multiclass concept) and
two sign crossings in $x$. Convolving with the Gaussian window of $Q_\mu=\mathcal N(\mu,1)$
and sweeping $\mu$ across $D$ yields $\Delta(\mu)$ with two sign changes, hence three
maximal sign-constant intervals in $\mu$. Identifiability of $\sgn\Delta$ over a
neighbourhood requires $\Delta$ to be sign-constant there, so the identifiable set has three
components, whereas the binary single-threshold benefit
$2\Phi(\theta-\mu)-\Phi(-\mu)-\Phi(a-\mu)$ is monotone in $\mu$ with a single flip.
\emph{(Validated: $\max|\Lambda|=0.70$; $\Lambda$ sign-crossings in $x=2$; $\Delta(\mu)$
sign-changes $=2$ ($3$ components) versus binary baseline $=1$; seed-robust;
\texttt{val\_multiclass\_capacity.py} Block~B.)}
\end{proof}

\noindent\emph{Relation to the dichotomy.} This is consistent with --- and sharper than ---
Lemma~\ref{lem:dich-multiflip} (which proved $\tau=1$ survives an \emph{observable}
multi-flip frontier as a local band-vs-frontier event). The point here is not merely that
the sign map flips often; it is that in the multi-conflict-pair regime the flip locus is
intrinsically a \emph{multi-component} object \emph{generated by the classifier geometry
itself}, even for canonical exponential-family concepts, so the alignment-blind, single-
direction margin that works in binary has no multiclass counterpart.

% ----------------------------------------------------------------------------
\subsection{Necessity II: the binary observable margin is the wrong reduct}
\label{sec:mc-thirdclass}
% ----------------------------------------------------------------------------
Even \emph{with} a single conflict pair, the standard label-free observable used in the
binary theory --- the deployed model's top score on $D$, i.e.\ the margin
$M=\mathbb E_{Q\mid D}[\text{top score}]-\tfrac12$ --- is no longer a sufficient statistic
for $\sgn\Delta$. The third-class residual confounds it.

\begin{theorem}[Third-class confounding counterexample; PROVEN, WITNESS]
\label{thm:mc-thirdclass}
Fix $Q_X=\mathcal N(0,1)$, $D=(0,a)$, and a single conflict pair $j_g=0$ (the deployed top
class), $j_f=1$. There exist two $K=3$ admissible target concepts $\eta^+,\eta^-$ with:
\begin{enumerate}\itemsep1pt
\item[\textnormal{(i)}] \textbf{identical} observable covariate law $Q_X$ (covariate shift:
$\eta$ does not affect $Q_X$), hence unlabeled $\TV=0$ for every statistic of any number of
unlabeled samples;
\item[\textnormal{(ii)}] \textbf{identical} deployed top-score profile $\eta_{0}(x)$ on $D$
(so the same binary-style margin $M$, the same predicted-class map, and the same top-label
calibration a label-free auditor can compute); and
\item[\textnormal{(iii)}] \textbf{opposite} benefit sign,
$\sgn\Delta(\eta^+)=+1$, $\sgn\Delta(\eta^-)=-1$.
\end{enumerate}
The distinguishing object is the split of the residual mass $1-\eta_0(x)$ between class
$1=j_f$ and the third class, which is invisible to $Q_X$ and to the top score. Consequently
a capacity built from the binary observable, $K_{\mathrm{bin}}=|M|/\beta$, can exceed $1$ in
\emph{both} worlds (predicting ``identifiable'') while $\sgn\Delta$ is \emph{not}
identifiable; the minimax label-free error is $\tfrac12$ (Le~Cam, $\TV=0$). Hence in
multiclass the binary margin $M$ is a \emph{strictly-too-coarse reduct}, and the admissible
budget must be placed as \textnormal{(M3)} on the contrast field $\Lambda$ rather than on the
top score.
\end{theorem}
\begin{proof}
Let $\eta_0(x)\equiv p$ on $D$ (a constant deployed top score, $p<\tfrac12$ so $M=p-\tfrac12<0$).
Put $\eta_1=(1-p)s(x)$, $\eta_2=(1-p)(1-s(x))$ with $s(x)\in[0,1]$ the residual split. Then
the benefit field is $\Lambda=\eta_1-\eta_0=(1-p)s-p$, and by \eqref{eq:mc-delta}
$\Delta=\int_D[(1-p)s(x)-p]\,\mathrm dQ$. Choosing $s\equiv s^+$ large makes $\Delta>0$ and
$s\equiv s^-$ small makes $\Delta<0$; both share $\eta_0\equiv p$ and the same $Q_X$. With
$p=0.45$, $s^+=0.95$, $s^-=0.05$, $a=2$: $\Delta^+=+0.035>0$ and $\Delta^-=-0.202<0$. The two
worlds have $\TV=0$ on every unlabeled statistic (the concept never enters $Q_X$), so Le~Cam's
two-point bound gives minimax error $\tfrac12$. Yet $|M|=0.05$, so $K_{\mathrm{bin}}=|M|/\beta>1$
for any $\beta<0.05$ in both worlds. \emph{(Validated: $\Delta^+=+0.0346$, $\Delta^-=-0.2016$,
opposite sign at shared $\eta_0=0.45$ and shared $Q_X$, with $K_{\mathrm{bin}}=2.5$ and $1.25$
at $\beta=0.02,0.04$ predicting ``identifiable'' against the truth;
\texttt{val\_multiclass\_capacity.py} Block~C.)}
\end{proof}

\noindent This matches the known unidentifiability of target accuracy from unlabeled data
alone absent constraints on the target conditional \citep{garg2022atc}, and the fact that the
natural multiclass calibration observable is the (weaker) top-label calibration
\citep{gupta2021toplabel}: the non-predicted-class contrast that drives $\Delta$ is exactly
the part those observables do not see.

% ----------------------------------------------------------------------------
\subsection{Honest status and the open core}\label{sec:mc-open}
% ----------------------------------------------------------------------------
\begin{remark}[What is and is not established]\label{rmk:mc-scope}
\textbf{Proven.} (a) The exact vector benefit identity \eqref{eq:mc-delta} for $K\ge3$
(Lemma~\ref{lem:mc-identity}). (b) A single-conflict-pair multiclass capacity at $\tau=1$
under the explicit hypotheses (M1)--(M3) (Theorem~\ref{thm:mc-singlepair}) --- a genuine
$K\ge3$ extension. (c) Two necessity results delimiting the wall: the conflict-pair
structure (M1) is necessary, since dropping it fragments the identifiable set into
$\ge3$ components even for canonical exponential-family concepts
(Theorem~\ref{thm:mc-breakerM1}); and the binary observable margin is a too-coarse reduct,
since two $K=3$ worlds share $(Q_X,\text{top score})$ yet have opposite $\sgn\Delta$ via the
third-class split (Theorem~\ref{thm:mc-thirdclass}). The distribution-free converse
(Theorem~\ref{thm:gen-df-conv}) already applies unchanged: whenever two admissible concepts
share $Q_X$ with opposite sign, minimax error is $\tfrac12$.

\textbf{Open.} A single \emph{computable} multiclass capacity \emph{without} (M1) --- the
genuinely multi-conflict-pair / vector-concept regime. Concretely: when $D$ carries several
conflict pairs and the admissible concept is a vector $\eta\in\Delta^{K-1}$ constrained by a
budget, the flip set is intrinsically multi-component (Theorem~\ref{thm:mc-breakerM1}) and
the worst-case flip over the admissible \emph{vector} class is not known to be a computable
functional of any natural label-free observable (Theorem~\ref{thm:mc-thirdclass} shows the
obvious candidate, the top-score margin, fails). By the dichotomy
(Theorem~\ref{thm:dich-main}) a capacity \emph{exists} iff the benefit sign factors through
the maximal reduct $O=[Q_X]$ \emph{computably}; the multiclass open problem is precisely
whether, and via which computable observable, that factorization can be \emph{exhibited} in
the vector-concept regime --- equivalently, what the weakest label-free supplement is that
pins the non-predicted-class contrast. We conjecture that, exactly as in binary
(Theorem~\ref{thm:conj1-dichotomy}), it is an irreducible supplement (here a per-conflict-
pair orientation on each component of the flip set, hence up to $\#\{\text{components}\}-1$
bits rather than one), but we do \emph{not} prove this.

\textbf{Verdict.} A strong \emph{partial} on the multiclass frontier: the binary capacity
lifts exactly to a single conflict pair and stops there for two named, witnessed reasons.
This converts the vague ``$K\ge3$ is open'' of Conjecture~\ref{conj:gen-capacity} into a
precise statement of \emph{what} lifts, \emph{what} is necessary, and \emph{what} the
residual open core is.
\end{remark}
